%! TeX root = ./main.tex
% vim: set spelllang=en:
\selectlanguage{english}
\twocolumn
\chapter*{Objects and notation}
\addcontentsline{toc}{chapter}{Objects and notation}
\markboth{\textsc{objects and notation}}{}

\paragraph{Graphs}
A graph~\(G\) is formed on a set of \emph{vertices}~\(\tnodes(G)\) connected by a set of
\emph{edges} or \emph{arcs} denoted~\(\tedges(G)\).
An edge connecting two vertices~\(a\) and~\(b\) is denoted~\(\{a, b\}\).
An arc starting at~\(a\) and ending at~\(b\) is denoted~\((a, b)\).

\paragraph{Trees}
The set of vertices of a tree~\(T\) is partitioned into the \emph{leaves} \(\tleavs(T)\) and the \emph{internal nodes} \(\tinner(T)\).
A \emph{rooted} tree has a root \(\troot(T) \in \tnodes(T)\) and is formed of arcs directed away from the root.
A tree \(T'\) \emph{displays} \(T\) if \(T\) is a minor of~\(T'\).
A \emph{forest} is a disjoint union of trees.

\paragraph{Paths}
The set of edges or arcs making up the shortest \emph{path} between two vertices \(v\) and \(w\) is denoted \(\tpath{v}{w}\).
The cardinality of this set defines the \emph{distance} between \(v\) and \(w\), denoted \(\tdist(v, w)\).

\vfill\break
\paragraph{Parenthood}
In a rooted tree \(T\), the parent of a vertex \(v\) other than \(\troot(T)\) is its unique predecessor, denoted \(\tparent(v)\).
The set of successors of an internal node \(v\), called its \emph{children}, is denoted \(\children(v)\).
These children are denoted \(v\lchild\) and \(v\rchild\) if there are two of them, or \(v\child\) if there is only one or if it is arbitrary.
The tree is \emph{binary} if the out-degree of all internal nodes is two, and \emph{partially binary} if it is one or two.

\paragraph{Descent}
If a vertex \(v\) is on the path leading from a root to a vertex \(w\), including~\(w\) itself, then \(v\) is an \emph{ancestor} of~\(w\), denoted \(v \tancs w\), and \(w\)~is a \emph{descendant} of \(v\), denoted \(w \tdesc v\).
When \(v\) and \(w\) are not comparable by \(\tdesc\), they are said to be \emph{separated}, denoted \(v \nottcomp w\).
The lowest common ancestor~(or \emph{lca}) of a set of vertices \(\{v_1, …, v_n\}\) is their ancestor furthest from the root, denoted \(\lca\{v_1, …, v_n\}\).
The subtree of \(T\) at \(v\), denoted \(\subtree{T}{v}\), is the subgraph of \(T\) formed by the descendants of \(v\).

% power set

\onecolumn
